% P19-E3: five original confidence experiments and Brier's original score.
\chapter{Calibration and Extreme Confidence}
\label{sec:p19-calibration-extreme-confidence}

A judgment can select the correct answer while misstating how reliably that
kind of judgment succeeds. Calibration asks whether reported probabilities
match outcome frequencies under a specified judgment--outcome law. It is
not synonymous with the proportion of answers correct, the informativeness
of the judgments, or their usefulness in a particular decision.

This chapter treats all five experiments in Fischhoff, Slovic and
Lichtenstein's \emph{Knowing with Certainty: The Appropriateness of Extreme
Confidence} \cite{FischhoffSlovicLichtenstein1977Certainty}, and Brier's
original three-page probability-forecast paper \cite{Brier1950Verification}.
The finite scoring and conditional-mean proofs below are handbook
derivations. The calibration decomposition is not attributed as a formula
appearing in Brier's 1950 article. Historical mortality rates and quiz answer
keys are identified as source task material, not current population facts.

\section{What is being calibrated?}
\label{sec:p19-calibration-law}

\begin{hbdefinition}{def:p19-calibration-law}{A joint law, accuracy and calibration}
Fix a finite probability space and jointly defined variables \(C,Y\), where
\(C\) takes values in a finite subset of \([0,1]\) and \(Y\in\{0,1\}\).
The judgment \(C\) reports the probability of the event \(Y=1\).
For each \(c\) with \(w_c=\mathbb P(C=c)>0\), put
\[
 m(c)=\mathbb E[Y\mid C=c]=\mathbb P(Y=1\mid C=c),\qquad
 \mu=\mathbb E[Y].
\]
The judgment is \emph{calibrated} under this law when \(m(c)=c\) on
every positive-mass judgment value. Assign \(m(c)=0\) on null judgment
classes to define \(m(C)\) everywhere. This arbitrary version value is
not a conditional probability inferred from an empty class.

For a quiz, \(Y=1\) may mean that the answer chosen by the participant is
correct under a specified answer key. Then \(\mu\) is mean answer accuracy.
For a judgment about a supplied statement, \(Y=1\) instead means that
statement is true; the participant need not have selected it. These are
different tasks and can induce different joint laws.

\emph{Mean overconfidence} is \(\mathbb E[C]-\mu>0\); conditional
overconfidence at \(c\) means \(c>m(c)\). Define the squared calibration
error and resolution by
\[
 \operatorname{Rel}=\mathbb E[(C-m(C))^2],\qquad
 \operatorname{Res}=\mathbb E[(m(C)-\mu)^2].
\]
Resolution measures variation in outcome rates between judgment classes.
It is not generally the variance of the reported probabilities themselves.
The binary squared-error score is \(L=\mathbb E[(C-Y)^2]\), with smaller
values better. These expectations use the finite conditional means of
\cref{def:p15-finite-conditional-expectation}.
\end{hbdefinition}

For observed pairs \((c_i,y_i)\), an empirical law assigns mass \(1/n\)
to each of \(n>0\) recorded trials. Its exact-value calibration table has
\(n_c\) reports equal to \(c\) and observed frequency
\(\widehat m(c)=\sum_{i:c_i=c}y_i/n_c\), for \(n_c>0\). Repeated trials from the same
person or the same question remain separate records in this descriptive law;
they are not thereby statistically independent observations from a population.
Coarser bins discard distinctions within them and can conceal opposing
calibration errors. A calibration claim must name its item, person and
circumstance mixture rather than refer to a person as calibrated everywhere.

\begin{hbtheorem}{thm:p19-calibration-decomposition}{proved=annex:p19.calibration-decomposition}{Conditional means, resolution and recalibration}
Under \cref{def:p19-calibration-law},
\[
 L=\operatorname{Rel}+\mathbb E[m(C)(1-m(C))]
  =\operatorname{Rel}+\mu(1-\mu)-\operatorname{Res}.
\]
For any real function \(g\) on the judgment values,
\[
 \mathbb E[(g(C)-Y)^2]
 =\mathbb E[(g(C)-m(C))^2]+\mathbb E[m(C)(1-m(C))].
\]
Consequently \(Q=m(C)\) is calibrated under the same law and reduces
the score by exactly \(\operatorname{Rel}\). It is the unique optimal
function of \(C\), up to null classes, for expected squared error.
Calibration implies \(\mathbb E[C]=\mu\), but equality of these means
does not imply calibration. Also
\(0\leq\operatorname{Res}\leq\mu(1-\mu)\).
\end{hbtheorem}

Using observed outcomes to fit \(\widehat m\) improves this empirical
score by the identity. It does not prove improved forecasting on new trials.
That would require a new-trial law and evidence about the fitted mapping's
behavior there. Likewise, the existence of an optimal conditional mean does
not establish that a participant knows or reports it.

\begin{hblemma}{thm:p19-binned-calibration}{proved=annex:p19.binned-calibration}{What a coarser calibration table omits}
Let \(B=h(C)\) be a finite bin label. On each positive-probability bin
define \(a_B=\mathbb E[C\mid B]\), \(b_B=\mathbb E[Y\mid B]\),
\(v_B=\mathbb E[(C-a_B)^2\mid B]\), and
\(\gamma_B=\mathbb E[(C-a_B)(Y-b_B)\mid B]\). These are random
variables constant within each bin. Then
\[
 L=\mathbb E[(a_B-b_B)^2]+\mu(1-\mu)
   -\mathbb E[(b_B-\mu)^2]+\mathbb E[v_B]-2\mathbb E[\gamma_B].
\]
Thus a three-term calibration decomposition based on arbitrary bins need
not equal the original score. The last two terms vanish when each bin
contains a single judgment value. Equality of bin means alone need not
establish calibration at the original values.
\end{hblemma}

\section{Brier's original scoring benchmark}
\label{sec:p19-brier-original}

Brier considered \(r\) mutually exclusive, exhaustive outcome categories.
On occasion \(i\), the probabilities \(f_{ij}\) sum to one; \(E_{ij}\)
is one for the category that occurred and zero otherwise. His score, denoted
\(P\) in the source and \(P_{\rm B}\) here, is
\[
 P_{\rm B}=\frac1n\sum_{i=1}^n\sum_{j=1}^r(f_{ij}-E_{ij})^2.
\]
The normalization divides by occasions, not by the number of categories
\cite[p.~1, eq.~(2)]{Brier1950Verification}.

\begin{hbtheorem}{thm:p19-brier-propriety}{proved=annex:p19.brier-propriety}{Category-sum score and strict propriety}
For \(r\geq2\), each category-sum loss lies in \([0,2]\). If the
outcome-category law is \(q=(q_1,\ldots,q_r)\), then a probability report
\(f\) has expected loss
\[
 \sum_j(f_j-q_j)^2+1-\sum_jq_j^2.
\]
It is uniquely minimized at \(f=q\); this is \emph{strict propriety} for
minimizing expected score. For a fixed empirical record, the best constant
forecast equals its category-frequency vector and scores
\(1-\sum_jq_j^2\).
For a binary forecast \((c,1-c)\) and outcome \((y,1-y)\), the
category-sum loss is \(2(c-y)^2\). Thus Brier's original binary scale is
twice \(L\), and its maximum is two rather than one.
\end{hbtheorem}

Strict propriety is a statement about optimization under a given probability
law and expected-score objective. Someone optimizing under an erroneous
belief can honestly report that belief and remain miscalibrated under the
law generating the outcomes. Monetary risk preferences can also differ from
minimizing expected score. The scoring rule itself is not evidence about
participants' calibration.

\begin{hbexample}{def:p19-brier-calibration-examples}{Actual forecast arithmetic and constructed comparisons}
Brier's ten-record example has rain probabilities
\[
 (.7,.9,.8,.4,.2,0,0,0,0,.1)
\]
and rain outcomes \((0,1,1,1,0,0,0,0,0,0)\). The original score is
.19; the one-coordinate score is .095. A constant .3 rain forecast scores
.42 on Brier's scale, and a constant .2 scores .44. The empirical .3
benchmark uses the outcome frequency in those ten records; it was not known
in advance \cite[p.~2, table~1]{Brier1950Verification}.

For a separate constructed comparison, let both classes have weight one
half. A constant forecast .75 with outcome rate .75 is calibrated, with
\(L=.1875\) and resolution zero. A forecast taking .6 and .9 with matching
conditional outcome rates is also calibrated, with resolution .0225 and
\(L=.165\). Interchanging the two conditional outcome rates gives the same
mean forecast and mean outcome .75, but calibration error .09 and
\(L=.255\). One bin covering both values would show equal mean judgment
and outcome and conceal the errors. These are stipulated laws, not extra
experimental observations.
If reports vary between .6 and .9 but both conditional outcome rates are
.75, their report variance is .0225 while resolution is zero; the loss
is .21. Variation in confidence alone does not supply resolution.
\end{hbexample}

\begin{hblemma}{thm:p19-score-calculations}{proved=annex:p19.score-calculations}{Scores, mean agreement and certainty thresholds}
The calculations in \cref{def:p19-brier-calibration-examples} hold. Odds
\(t:1\), for \(t\geq0\), correspond to probability \(t/(1+t)\).
Every finite \(t\) gives a probability strictly below one. Under a calibrated
law, if \(A_t=\{C\geq t/(1+t)\}\) has positive probability, then
\[
 \mathbb P(Y=0\mid A_t)\leq\frac1{1+t}.
\]
Furthermore \(\mathbb P(C=1,Y=0)=\mathbb P(C=0,Y=1)=0\).
These are properties of the specified joint law; a finite observed error
frequency at a finite odds threshold is not required to equal its expectation.
\end{hblemma}

Brier also tabulated 85 forecasts in five bins. The observed rain proportions
were .07, .10, .29, .40 and .50 for reported intervals .00--.19,
.20--.39, .40--.59, .60--.79 and .80--1.00, respectively. The paper
does not give their individual forecasts or numbers per bin, so those rows
do not determine an overall score, standard error or complete calibration
law. Brier reported using the score over several months for experimental
six-hour precipitation forecasts by the Weather Bureau's Short Range
Forecast Development Section, while noting questions needing further study
\cite[pp.~2--3, table~2]{Brier1950Verification}.

The decomposition above states how calibration and resolution affect squared
error. Improving the score need not increase the frequency of extreme reports
or a correlation coefficient. For example, moving a constant report .9 to
.6 when the outcome rate is .6 lowers binary loss from .33 to .24 while
making the report less extreme. Such constant forecasts have no nonzero
forecast variance from which to form an ordinary correlation coefficient.

\section{Experiment 1: extreme probability reports in four formats}
\label{sec:p19-certainty-experiment-one}

Fischhoff, Slovic and Lichtenstein's first experiment used general-knowledge
questions about history, music, geography, nature and literature. The four
formats were: complete an open-ended question and assess the answer on
\([0,1]\); judge a supplied single statement on \([0,1]\); choose between
two alternatives and assess the chosen answer on \([.5,1]\); or judge
an experimenter-selected alternative from a displayed pair on \([0,1]\).
In the last format the experimenter selected the alternative randomly.
The experimenters assigned participants to formats according to preferred
session time/date, not reported random assignment. Items also differed
somewhat between formats \cite[pp.~553--554]{FischhoffSlovicLichtenstein1977Certainty}.

The method reports 361 paid volunteers recruited through an advertisement
in the University of Oregon student newspaper. Table~1 gives the following
condition fields. For a report of zero, the last column is the proportion
of judged statements actually true, which would be zero under perfect
certainty calibration; it is not a proportion of correct zero reports.
\begin{center}\small
\begin{tabular}{lrrrrrr}
Format&Items&People&Responses&Report&\% at report&\% true\\\hline
Open answer&43&30&1290&1&19.7&83.1\\
One statement&75&86&6450&1&14.2&71.7\\
&&&&0&13.8&29.5\\
Choose one of two&75&120&9000&1&21.8&81.8\\
Supplied one of two&50&131&6500&1&17.3&80.7\\
&&&&0&19.1&20.5
\end{tabular}
\end{center}
These printed fields are not fully reconcilable: the people column sums to
367 rather than the prose's 361, and \(131\times50=6550\), not 6500.
No unreported exclusion is invented to resolve them. The last two columns
are source percentages, not counts reconstructed from rounded values.

Certainty-one reports were correct in 71.7--83.1\% of cases, hence wrong
in 16.9--28.3\%. Certainty-zero reports were true in 20.5--29.5\%.
These are the table-supported directions of error.
For the full-range judgments, a report of zero expresses certainty in the
statement's negation. Errors at both endpoints therefore need not yield a
positive mean bias in the reported probability of the statement itself.
Comparisons among formats
do not isolate a format effect because both assignment and item sets differed.
The article examines extreme reports here; its introductory Figure~1 summarizes
earlier calibration work rather than supplying a sixth new experiment.

\section{Experiments 2 and 3: odds and instruction}
\label{sec:p19-certainty-odds-experiments}

Experiment~2 used 66 paid volunteers from the same newspaper-advertisement
recruitment route. Each judged which of two causes of death was more frequent
in the United States and then reported odds that the chosen answer was
correct. The 106 pairs came from 41 causes, with each cause in approximately
six pairs. The authors describe rate ratios from about 1.25:1 to 100000:1.
Their rate table also includes zero for smallpox, so a positive-to-zero
ratio cannot be included in a finite range. The complete pair list is not
printed, and this range is not assumed to cover every possible pairing.
A glossary
was supplied for ambiguous cause descriptions. Source rates came from vital
statistics and Metropolitan Life Insurance reports; the authors describe them
as fairly stable over the preceding five years of available statistics
\cite[pp.~554--556]{FischhoffSlovicLichtenstein1977Certainty}.

For completeness, the historical task rates per year per 100 million United
States residents were as follows. They are answer-key inputs for this study,
not updated mortality estimates.
\begin{center}\small
\begin{tabular}{p{.30\linewidth}r|p{.30\linewidth}r}
Cause&Rate&Cause&Rate\\\hline
Smallpox&0&Firearm accident&1100\\
Vitamin poisoning&.5&Solid/liquid poisoning&1250\\
Botulism&1&Tuberculosis&1800\\
Measles&2.4&Fire/flames&3600\\
Fireworks&3&Drowning&3600\\
Smallpox vaccination&4&Leukemia&7100\\
Whooping cough&7.2&Accidental falls&8500\\
Polio&8.3&Homicide&9200\\
Venomous bite/sting&23.5&Emphysema&10600\\
Tornado&44&Suicide&12000\\
Lightning&52&Breast cancer&15200\\
Nonvenomous animal&63&Diabetes&19000\\
Flood&100&Motor-vehicle accident&27000\\
Excess cold&163&Lung cancer&37000\\
Syphilis&200&Digestive-system cancer&46400\\
Pregnancy, childbirth, abortion&220&All accidents&55000\\
Infectious hepatitis&330&Stroke&102000\\
Appendicitis&440&All cancers&160000\\
Electrocution&500&Heart disease&360000\\
Motor-vehicle/train collision&740&All diseases&849000\\
Asthma&920&&
\end{tabular}
\end{center}

The odds instruction explained equal odds as guessing and permitted any
value at least 1:1, including values between marked powers of ten and values
above a million to one. A report below 1:1 would indicate that the other
answer should have been chosen. Thus the open-ended odds scale avoided
restricting high confidence to a few probability ticks such as .90, .95 and
1.00. There were 6996 judgments, equal to \(66\times106\).

Experiment~3 used the same 106 questions and odds format with 40 paid
advertisement respondents, in groups of about twenty. It began with a
twenty-minute lecture explaining probability, odds, low-odds responses,
their conversion chart and calibration. Participants were encouraged to ask
questions and to aim for calibration. The full lecture text is not printed
in the article \cite[p.~557]{FischhoffSlovicLichtenstein1977Certainty}.

The following retains the counts and percent correct in the major odds
categories from Table~3 for Experiments~2, 3 and~4. The final category is
discussed in the text as a million to one or higher. Intermediate reports
are not all included in these rows.
\begin{center}\small
\begin{tabular}{rrrrrrr}
Odds to one&E2 count&E2 \% correct&E3 count&E3 \% correct&E4 count&E4 \% correct\\\hline
1&644&53&339&54&861&53\\
1.5&68&57&108&59&210&56\\
2&575&64&434&65&455&63\\
3&189&71&252&65&157&76\\
5&250&70&322&71&194&76\\
10&1167&66&390&76&376&74\\
20&126&72&163&81&66&85\\
50&258&68&227&74&69&83\\
100&1180&73&319&87&376&80\\
1000&862&81&219&84&334&88\\
10000&459&87&138&92&263&89\\
100000&163&85&23&96&134&92\\
1000000&157&90&47&96&360&94\\\hline
Listed total&6098&&2981&&3855&
\end{tabular}
\end{center}
For comparison, calibrated probabilities at odds 1, 1.5, 2, 3, 5, 10,
20, 50, 100 and 1000 to one are respectively 50\%, 60\%,
\(66\tfrac23\%\), 75\%, \(83\tfrac13\%\), about 90.91\%,
95.24\%, 98.04\%, 99.01\% and 99.90\%. At still larger finite
odds they approach but do not equal 100\%.

Experiment~2 had 3560 reports above 50:1, about 51\% of all judgments;
almost a quarter exceeded 1000:1. At 100:1 only 73\% were correct;
the million-or-higher group was 90\% correct. Higher confidence therefore
did not track the near-certainty implied by those numerical odds.

Experiment~3 used more intermediate odds and fewer large odds. The source
reports about 43\% at 5:1 or below, against about 27\% in the earlier
odds experiment, and about a third at 50:1 or above, of which approximately
83\% were correct. The high-odds entries were generally better calibrated
than Experiment~2's, but remained far from their benchmarks. This is a
comparison of separate cohorts; no randomized instruction-only comparison
is reported.

Some printed category shares conflict with the counts. The listed totals
6098, 2981 and 3855 are respectively about 87.2\%, 70.3\% and 86.6\%
of the reported or nominal complete trial totals 6996, 4240 and 4452.
The table instead prints 88\%, 70\% and 75\%; its Experiment~4 row
for odds 2:1 prints 1\% despite 455 records. The counts and percent-correct
columns above are preserved without adopting the inconsistent shares.
Experiment~3's omitted share is about 29.7\% of its nominal trials;
the prose's reference to a remaining 32\% is not used as a denominator.

\section{Experiments 4 and 5: general knowledge and betting}
\label{sec:p19-certainty-betting-experiments}

Experiment~4 recruited 42 paid volunteers through the Oregon student
newspaper. It used 106 two-alternative general-knowledge questions selected
from a characterized item pool and matched individually in difficulty to
the lethal-event questions. The instruction lecture paralleled Experiment~3.
Participants selected answers and reported odds. Table~3 reports overall
percent correct as 71.0, 72.5 and 73.1 for Experiments~2, 3 and~4.
These overall accuracy figures are different from calibration within an
odds category. Experiment~4 included examples about magazine circulation,
colonial history and bile pigments; about one quarter of reports were at
least 1000:1 \cite[pp.~557--558]{FischhoffSlovicLichtenstein1977Certainty}.

After the questions, the experimenter offered a game using answers assigned
odds of at least 50:1. For each such answer a chip would be drawn from a
bag with 100 white and 2 red chips, replacing it before the next draw.
A wrong answer would cost the participant one dollar; each red chip would
pay the participant one dollar. The announced participation payment was
\$2.50. Those declining the game were successively offered \$1.50, \$2
and \$2.50 per red chip, while their wrong-answer cost stayed one dollar.
The game was hypothetical in Experiment~4; no money was wagered.

The paper reports 27 acceptances at \$1, then six at \$1.50, three at
\$2 and two at \$2.50, together with three refusals. Those incremental
acceptances total 38, whereas the subsequent analysis uses 39 acceptors
and the refusal statement would also imply 39. The discrepancy is left
explicit; one participant's stake is not supplied by assumption.
No participant said they would change their decision when asked whether
the game should be played immediately for real money. Two requested the
opportunity, which the experimenter declined.

The reported simulation for 39 acceptors used 1495 high-odds answers,
averaging 38.3 per person. Thirty-six had negative expected net money and
three positive. Expected results ranged from a loss of \$25.63 to a gain
of \$1.84; mean loss was \$3.64, median loss \$2.35, and ten people
had expected losses above \$5. The total expected loss was \$142.13,
about 9.5 cents per selected answer. The two people who particularly
wanted to play had expected losses totaling \$33.38. These are simulated
expected payments, not realized losses
\cite[pp.~558--559]{FischhoffSlovicLichtenstein1977Certainty}.

Experiment~5 involved 19 participants and otherwise followed Experiment~4,
but announced the game as real. Six accepted at \$1, three more at
\$1.50, one at \$2 and three at \$2.50; six declined, some citing
miscalibration and others dislike of gambling. Participants initially were
told they could lose their participation pay and more. After they chose
whether to play, they were told that winnings would be added to their pay
but losses would not reduce the originally promised payment. The game was
played under that protection. Thus willingness was elicited before the
protection announcement, while realized payments used the protected rule.

The 13 players missed 46 of their 387 selected answers, 11.9\%. No red
chip was drawn. Their unprotected game balances would all have been negative,
from \$1 to \$11; the paper does not report confiscating these amounts.
The simulated expected unprotected loss averaged \$2.64 per player, with
four above \$6. The result concerns the 13 self-selected players, not a
complete calibration curve for all nineteen participants
\cite[p.~559]{FischhoffSlovicLichtenstein1977Certainty}.

\begin{hblemma}{thm:p19-certainty-game}{proved=annex:p19.certainty-game}{Expected money in the chip comparison}
For \(k\) selected answers with \(w\) known wrong answers, let each
draw have red probability \(1/51\), and let \(b\) dollars be paid per
red chip. The expected unprotected net payment to the participant is
\[
 \frac{kb}{51}-w.
\]
This identity needs the marginal probability of each draw, not independence
between answers. If a participant's probability for correctness of answer
\(i\), conditional on their information, is \(q_i\geq50/51\), and they
assign each chip draw red probability \(1/51\), their
subjective expected unprotected net is
\(kb/51-\sum_i(1-q_i)\geq k(b-1)/51\).
At \(b=1\) this is nonnegative, and strictly positive only if some
\(q_i>50/51\). A score or expected-money inequality alone does not
require acceptance by every person with arbitrary risk preferences.
\end{hblemma}

For the source's example of seventeen selected answers at \(b=1\),
expected chip payment is \(17/51=1/3\) dollar before wrong-answer losses.
A bag with 14 white and 2 red chips would have red probability \(1/8\),
matching an error rate of one eighth at equal stakes. These expectations
explain the experimenters' comparison; the game is not a proper probability-
report scoring rule administered prospectively to elicit every confidence
judgment. Refusing it can reflect risk preferences as well as beliefs.

\section{How widely were extreme errors distributed?}
\label{sec:p19-certainty-item-person-analyses}

The source defines an extreme error as a wrong answer given odds at least
50:1. Its Table~4 gives the complete following distributions over people.
The more-than-seventeen category gives each exceptional person's actual
number of errors in parentheses.
\begin{center}\small
\begin{tabular}{r|rr@{\qquad}r|rr}
Errors&E3 people&E4 people&Errors&E3 people&E4 people\\\hline
0&5&3&10&2&2\\
1&3&7&11&0&1\\
2&6&5&12&1&0\\
3&6&5&13&1&0\\
4&4&9&14&0&0\\
5&1&1&15&0&2\\
6&4&1&16&2&0\\
7&2&2&17&1&0\\
8&1&0&More than 17&1 (32)&1 (27)\\
9&0&3&&&
\end{tabular}
\end{center}
The corresponding item counts, ordered by the number of people making an
extreme error on an item from zero through fifteen, were
\[
\begin{array}{c|rrrrrrrrrrrrrrrr}
 &0&1&2&3&4&5&6&7&8&9&10&11&12&13&14&15\\\hline
E3&33&25&19&13&3&2&4&1&1&0&0&2&1&0&1&1\\
E4&41&20&16&10&8&2&1&3&0&0&3&0&0&1&0&1
\end{array}
\]
Thus five of forty and three of forty-two people had no extreme error in
these tasks. Thirty-three and forty-one of the respective 106 items had
none. The median numbers of extreme errors per person were 3.5 and 4.
The authors report that deleting the people with 32 and 27 such errors
did not change their conclusions, without publishing replacement tables
\cite[pp.~560--561]{FischhoffSlovicLichtenstein1977Certainty}.

An item was called \emph{deceptive} when fewer than half the participants
answered correctly. There were eighteen such items in Experiment~3 and
seventeen in Experiment~4. The following percentages wrong at successively
higher reported-odds thresholds reproduce Table~5; the parentheses are
numbers of items, not response denominators.
\begin{center}\small
\begin{tabular}{lrrr}
Item set&\(\geq50:1\)&\(\geq100:1\)&\(\geq1000:1\)\\\hline
E3 all (106)&16.6&14.1&12.9\\
E3 deceptive (18)&73.9&75.5&72.3\\
E3 other (88)&8.9&6.7&6.9\\
E4 all (106)&13.8&13.1&10.8\\
E4 deceptive (17)&73.2&76.7&70.6\\
E4 other (89)&7.6&6.8&5.2
\end{tabular}
\end{center}
Under a calibrated joint law, the exact corresponding upper bounds on the
\emph{expected} error percentages are \(100/51\), \(100/101\) and
\(100/1001\). Rounded figures such as 1.96 and .99 must not replace
these exact bounds in a strict inequality. Response counts for every
threshold-by-item subset are not printed, and the subsets are nested and
repeated within people. These percentages do not supply independent samples
for a new significance calculation.

The source reports that one third of items answered correctly by at least
90\% of participants still had a wrong answer assigned at least 1000:1,
and that deleting the two extreme participants changed this little. It
describes the nondeceptive-item error-count distribution as resembling a
Poisson distribution; it supplies no fitted parameter, goodness-of-fit test
or proof of independent Poisson errors.

The six lethal-event examples in Table~6 were pregnancy/appendicitis,
all accidents/stroke, homicide/suicide, measles/fireworks, suicide/diabetes
and breast cancer/diabetes. In each pair, the first historical rate above
is lower. Their respective percent-correct figures were
15, 17.5, 25, 25, 27.5 and 30; numbers of extreme errors were
15, 14, 12, 5, 8 and 1 out of forty participants.

Table~7 supplies seven general-knowledge examples. Listed here as source
answer-key entries, with percent correct and number of extreme errors out
of forty-two, they are: cacao production, Africa rather than South America
(4.8; 15); mortality, appendicitis rather than pregnancy/abortion/childbirth
(19.0; 13); first air raid, 1849 rather than 1937 (26.2; 10); Adonis,
vegetation rather than love (31.0; 10); Gibran's principal religious
inspiration, Christian rather than Buddhist (33.3; 7); \emph{Dido and
Aeneas}, Purcell rather than Berlioz (33.3; 2); and potato origin, Peru
rather than Ireland (35.7; 10). Some stems were abbreviated in the source.
Their interpretation depends on that historical item wording and answer key,
not a claim that every abbreviated cultural question has one uncontested
answer in all contexts \cite[p.~562]{FischhoffSlovicLichtenstein1977Certainty}.

\section{Evidence limits and explanatory alternatives}
\label{sec:p19-certainty-evidence-limits}

The studies show substantial extreme-error frequencies under their stated
tasks and answer keys. They do not establish that every individual is
miscalibrated in every setting. Items and volunteers were not probability
samples of everyday reasoning; item difficulty was deliberately varied and
matched. The authors explicitly identify representativeness of difficulty
as a central unresolved generalization question.

Their proposed explanations include reliance on salient remembered examples,
unexamined assumptions in inference chains, loss of uncertainty between
inferential stages, and reconstructive memory. They call these proposals
a framework for further research rather than an established explanation
of the results. Detailed process protocols were absent. Experiment~3
participants were asked to select one answer they considered certain and
write why they were confident; the paper prints one such explanation in
note~4. This limited report is not a complete record of every inference.
Aggregate calibration curves therefore do not identify one mechanism
\cite[pp.~561--564]{FischhoffSlovicLichtenstein1977Certainty}.

The authors also report that encyclopedia and dictionary answer sources
sometimes disagreed. They removed offending items before the experiments
and checked the remaining answers again; the number removed is not given.
No further participant-exclusion rule is reported. Demographic distributions,
recruitment dates, complete question lists, the full glossary and lecture,
participant-level records and confidence intervals are not supplied here.
The reported comparisons are descriptive; no inferential tests for the
principal experiment contrasts are printed. Between-experiment differences
do not isolate instruction, content or incentives through a randomized
factorial design. These limits coexist with the reported extreme errors;
they are not resolved by choosing a proper score or a larger apparent
number of repeated judgments.
